% !TEX program = pdflatex
\documentclass[10pt]{article}
\usepackage[utf8]{inputenc}
\usepackage[T1]{fontenc}
\usepackage[english]{babel}
\usepackage{lmodern,microtype}
\usepackage[margin=0.72in]{geometry}
\usepackage{amsmath,amssymb,amsthm,mathtools}
\usepackage{booktabs,array,xcolor,enumitem}
\usepackage[most]{tcolorbox}
\usepackage{fancyhdr}
\usepackage{needspace}
\usepackage[hidelinks]{hyperref}

\definecolor{CourseBlue}{HTML}{1A3C68}
\definecolor{CourseGreen}{HTML}{216E39}
\definecolor{CourseGold}{HTML}{B8860B}
\definecolor{CourseRed}{HTML}{9A2C2C}
\definecolor{SoftBlue}{HTML}{F2F6FA}
\definecolor{SoftGreen}{HTML}{F1F8F3}
\definecolor{SoftGold}{HTML}{FFF8E8}
\definecolor{SoftRed}{HTML}{FFF3F3}
\setlength{\parindent}{0pt}
\setlength{\parskip}{0.32em}
\setlist[itemize]{leftmargin=1.5em,itemsep=0.12em,topsep=0.2em}
\setlist[enumerate]{leftmargin=1.7em,itemsep=0.12em,topsep=0.2em}
\allowdisplaybreaks

\pagestyle{fancy}
\fancyhf{}
\lhead{Stochastic Calculus and Option Pricing}
\rhead{Lecture 10 Notes}
\cfoot{\thepage}
\renewcommand{\headrulewidth}{0.3pt}

\newcommand{\E}{\mathbb E}
\newcommand{\Pp}{\mathbb P}
\newcommand{\Q}{\mathbb Q}
\newcommand{\R}{\mathbb R}
\newcommand{\1}{\mathbf 1}
\newcommand{\cF}{\mathcal F}
\newcommand{\dd}{\,\mathrm d}
\newcommand{\esssup}{\operatorname*{ess\,sup}}
\theoremstyle{definition}
\newtheorem{definition}{Definition}[section]
\newtheorem{proposition}[definition]{Proposition}
\newtheorem{theorem}[definition]{Theorem}
\newtheorem{remark}[definition]{Remark}

\newtcolorbox{boardexample}[1]{enhanced,breakable,colback=SoftBlue,colframe=CourseBlue,boxrule=0.55pt,arc=1mm,left=1.5mm,right=1.5mm,top=1mm,bottom=1mm,before skip=0.55em,after skip=0.55em,title={Example --- #1},fonttitle=\bfseries}
\newtcolorbox{studentexercise}[1]{enhanced,breakable,colback=SoftGold,colframe=CourseGold,boxrule=0.55pt,arc=1mm,left=1.5mm,right=1.5mm,top=1mm,bottom=1mm,before skip=0.55em,after skip=0.35em,title={Exercise --- #1},fonttitle=\bfseries}
\newtcolorbox{shortsolution}[1][]{enhanced,breakable,colback=SoftGreen,colframe=CourseGreen,boxrule=0.55pt,arc=1mm,left=1.5mm,right=1.5mm,top=1mm,bottom=1mm,before skip=0.25em,after skip=0.6em,title={Short solution#1},fonttitle=\bfseries}
\newtcolorbox{keyidea}[1]{enhanced,breakable,colback=white,colframe=CourseGreen,boxrule=0.7pt,arc=1mm,left=1.5mm,right=1.5mm,top=1mm,bottom=1mm,before skip=0.5em,after skip=0.5em,title={#1},fonttitle=\bfseries}
\newtcolorbox{warningbox}[1]{enhanced,breakable,colback=SoftRed,colframe=CourseRed,boxrule=0.65pt,arc=1mm,left=1.5mm,right=1mm,top=1mm,bottom=1mm,before skip=0.5em,after skip=0.5em,title={#1},fonttitle=\bfseries}
\newtcolorbox{assumptionbox}[1]{enhanced,breakable,colback=SoftGold,colframe=CourseGold,boxrule=0.65pt,arc=1mm,left=1.5mm,right=1.5mm,top=1mm,bottom=1mm,before skip=0.5em,after skip=0.5em,title={#1},fonttitle=\bfseries}

\title{\textbf{American Options and Optimal Stopping}\\
\large Lecture 10 Notes --- Exercise Rights, Dynamic Programming, and Free Boundaries}
\author{}\date{}
\begin{document}
\maketitle
\vspace{-1.4em}

\begin{keyidea}{Goals and logical structure}
European claims have one prescribed cash-flow date. An American claim lets
its holder choose a stopping time, so its price combines no-arbitrage
valuation with a timing decision:
\[
\begin{aligned}
\text{exercise right}&\longrightarrow\text{optimal stopping under }\Q\\
&\longrightarrow\text{backward induction / Snell envelope}\\
&\longrightarrow\text{obstacle PDE and numerical methods}.
\end{aligned}
\]
At every date ask one question: is immediate exercise worth more than
continuation?
\end{keyidea}

% ============================================================
\section{The exercise right changes the pricing problem}
% ============================================================

\subsection{European and American exercise}

\begin{definition}[European and American claims]
Let \(g_t\) be the cash payoff received if a claim is exercised at date
\(t\). A \emph{European} claim pays \(g_T\) at one fixed maturity \(T\).
An \emph{American} claim may be exercised at any stopping time
\[
\tau\in\mathcal T_{t,T}
=\{\tau:\tau\text{ is a stopping time and }t\leq\tau\leq T\}.
\]
Its cash flow is \(g_\tau\).
\end{definition}

For a stock \(S\) and strike \(K\), the exercise payoffs are
\[
g^{\rm call}(S_t)=(S_t-K)^+,\qquad
g^{\rm put}(S_t)=(K-S_t)^+.
\]
The word American describes the exercise convention. It does not describe
the stock or the exchange.

\begin{boardexample}{more rights cannot reduce value}
The holder of an American put can choose \(\tau=T\), reproducing the
European exercise policy. Hence
\[
V_t^{\rm Am}\geq V_t^{\rm Eur}.
\]
The difference is the \emph{early-exercise premium}; it can be zero even
though the American holder owns an additional right.
\end{boardexample}

\subsection{The seller's super-hedging viewpoint}

Let \(B_t=e^{rt}\) be the bank account and \(X^\pi\) the wealth of a
self-financing strategy. The writer must be ready whenever the holder
exercises. Thus \(\pi\) is a super-hedge when
\[
X^\pi_u\geq g_u\quad\text{for all }u\in[t,T],
\]
equivalently \(X^\pi_\tau\geq g_\tau\) for every admissible stopping time.

\begin{keyidea}{Two equivalent views in a complete market}
\[
\boxed{
V_t^{\rm Am}
=\inf_{\pi}\{X^\pi_t:X^\pi_u\geq g_u\ \forall u\in[t,T]\}
=\esssup_{\tau\in\mathcal T_{t,T}}
\E^\Q[e^{-r(\tau-t)}g_\tau\mid\cF_t].}
\]
The first formula protects the writer; the second describes the holder's
optimal decision. Completeness is the bridge between them.
\end{keyidea}

\begin{warningbox}{A stopping time is essential}
The holder cannot select an exercise date after seeing future prices. The
exercise decision at date \(t\) may use only \(\cF_t\). A future-dependent
choice would overstate the contract value.
\end{warningbox}

\begin{studentexercise}{first comparison}
Which inequality always holds between American and European options with the
same payoff and maturity? Does it mean that exercise is optimal at time zero?
\end{studentexercise}
\begin{shortsolution}
\(V_0^{\rm Am}\geq V_0^{\rm Eur}\), because waiting to maturity is
available. It does not identify the optimal policy: the holder may continue
at every date before maturity, in which case the two values coincide.
\end{shortsolution}

% ============================================================
\section{Discrete time: backward induction}
% ============================================================

\subsection{Risk-neutral recursion}

Consider dates \(0,1,\ldots,N\), a bank account growing by \(R=1+r>0\),
and an equivalent martingale measure \(\Q\). Let \(Y_n=g_n(S_n)\). The
American value is
\[
V_n=\esssup_{\tau\in\mathcal T_{n,N}}
\E^\Q[R^{-(\tau-n)}Y_\tau\mid\cF_n].
\]
At maturity there is no remaining choice: \(V_N=Y_N\).

\medskip
\begin{theorem}[Bellman recursion]
\[
\boxed{
V_n=\max\left\{Y_n,\frac1R\E^\Q[V_{n+1}\mid\cF_n]\right\},
\qquad V_N=Y_N.}
\]
\end{theorem}
\begin{proof}
Any exercise policy either stops now, giving \(Y_n\), or continues for at
least one period. In the latter case the best remaining value is
\(R^{-1}\E^\Q[V_{n+1}\mid\cF_n]\). Take the better alternative.
\end{proof}

\begin{assumptionbox}{Exercise versus continuation}
Define \(E_n=Y_n\) and
\[
C_n=\frac1R\E^\Q[V_{n+1}\mid\cF_n].
\]
Then \(E_n>C_n\) means exercise, \(E_n<C_n\) means continue, and equality
means that either action is optimal.
\end{assumptionbox}

\subsection{A binomial implementation}

In a binomial model
\[
S_{n+1}=\begin{cases}uS_n,&\text{up},\\dS_n,&\text{down},\end{cases}
\qquad d<R<u,
\]
the risk-neutral up probability is \(p=(R-d)/(u-d)\). Thus
\[
C_n=\frac{pV_{n+1}^{u}+(1-p)V_{n+1}^{d}}{R},
\qquad V_n=\max\{g(S_n),C_n\}.
\]
Work from terminal payoffs backward through the tree, recording nodes where
the intrinsic value wins.

\begin{boardexample}{a strict early-exercise node}
Take \(S_0=100\), \(u=1.2\), \(d=0.8\), \(R=1.05\), and a put with \(K=105\).
Then \(p=0.625\). At a date-1 down node \(S_1=80\), exercise pays \(25\).
The two terminal put values are \(9\) and \(41\), so
\[
C_1(80)=\frac{0.625(9)+0.375(41)}{1.05}=20.
\]
Since \(25>20\), early exercise is strictly optimal at this node.
\end{boardexample}

The up-node value is \(V_1(120)=3.214\), hence
\[
V_0^{\rm Am}
=\frac{0.625(3.214)+0.375(25)}{1.05}\approx10.84.
\]
The European value is approximately \(9.06\), leaving an early-exercise
premium of about \(1.78\).

\begin{studentexercise}{one backward step}
Let \(R=1.02\), \(p=0.6\), and \(K=100\). At a node with \(S=85\), an
American put has terminal continuation values \(4\) and \(22\). Exercise,
continue, or indifferent?
\end{studentexercise}
\begin{shortsolution}
Exercise pays \(15\), while continuation is
\[
\frac{0.6(4)+0.4(22)}{1.02}\approx10.98.
\]
The holder exercises. Use terminal option values, not terminal stock prices,
in the continuation calculation.
\end{shortsolution}

\subsection{Snell envelope and the hedge}

\begin{definition}[Discounted Snell envelope]
With \(Z_n=R^{-n}Y_n\), define
\[
U_n=\esssup_{\tau\in\mathcal T_{n,N}}\E^\Q[Z_\tau\mid\cF_n].
\]
\(U\) is the smallest \(\Q\)-supermartingale dominating \(Z\), and
\(V_n=R^nU_n\).
\end{definition}

\begin{keyidea}{Optimal policy}
One optimal stopping rule is the first entry into the exercise region:
\[
\tau_n^\star=\inf\{k\in\{n,\ldots,N\}:V_k=Y_k\}.
\]
In a complete market, the martingale part of the Snell envelope yields the
dynamic hedge for the seller.
\end{keyidea}

% ============================================================
\section{Continuous time: stopping regions and boundaries}
% ============================================================

\subsection{Black--Scholes formulation}

Under \(\Q\), suppose
\[
\frac{\dd S_t}{S_t}=(r-q)\dd t+\sigma\dd W_t^\Q,
\qquad\sigma>0.
\]
For exercise payoff \(g(S_t)\), the Markov value is
\[
\boxed{
v(t,s)=\sup_{\tau\in\mathcal T_{t,T}}
\E^\Q_{t,s}[e^{-r(\tau-t)}g(S_\tau)].}
\]
The continuation and exercise regions are respectively where
\(v(t,s)>g(s)\) and \(v(t,s)=g(s)\).

\subsection{A non-dividend-paying call}

\begin{proposition}[No early exercise for a standard call]
If \(q=0\) and \(r\geq0\), then
\[
C^{\rm Am}(t,s)=C^{\rm Eur}(t,s)
\]
for a call payoff \((s-K)^+\). Early exercise is never strictly optimal
before maturity.
\end{proposition}
\begin{proof}
For \(\theta=T-t>0\),
\[
C^{\rm Eur}(t,s)
\geq e^{-r\theta}\E^\Q_{t,s}[S_T-K]
=s-Ke^{-r\theta}\geq s-K.
\]
Together with nonnegativity, this gives
\(C^{\rm Eur}(t,s)\geq(s-K)^+\). At any stopping time, retaining the
European call is at least as valuable as exercising it, and its discounted
price is a martingale.
\end{proof}

\begin{warningbox}{The dividend qualification matters}
When \(q>0\), holding the stock before an ex-dividend date can be valuable.
An American call may then be exercised early. Never quote the preceding
result without the no-dividend assumption.
\end{warningbox}

\subsection{American puts and boundaries}

For a put, receiving \(K-S_t\) now can be attractive at low stock values.
In the standard model with \(r>0\), the exercise region typically has the
form
\[
\mathcal E=\{(t,s):s\leq b(t)\},\qquad
\mathcal C=\{(t,s):s>b(t)\},
\]
where \(b(t)<K\) before maturity. Deep below \(K\), waiting delays cash while
adding little useful optionality; near the strike, continuation is valuable.

% ============================================================
\section{The obstacle PDE}
% ============================================================

\subsection{Variational inequality and free boundary}

Let
\[
\mathcal Lf=(r-q)sf_s+\frac12\sigma^2s^2f_{ss}.
\]
In the continuation region, hedging gives
\[
v_t+\mathcal Lv-rv=0.
\]
Together with the requirement \(v\geq g\), this becomes
\[
\boxed{
\min\{v-g,-(v_t+\mathcal Lv-rv)\}=0,\qquad v(T,s)=g(s).}
\]
For a regular American put boundary \(s=b(t)\), value matching and smooth
fit read
\[
v(t,b(t))=K-b(t),\qquad v_s(t,b(t))=-1.
\]
The boundary is free: it is a result of the pricing problem, not a contract
input.

\subsection{A complete PDE problem for the American put}

For the numerical example below, specialize to no dividends \(q=0\),
\(r\geq0\), and \(g(s)=(K-s)^+\). Define the pricing operator,
including discounting, by
\[
\mathcal Af=rsf_s+\tfrac12\sigma^2s^2f_{ss}-rf.
\]
The variational inequality is equivalent, wherever the derivatives exist, to
\[
\boxed{
v-g\geq0,\qquad -(v_t+\mathcal Av)\geq0,\qquad
(v-g)(v_t+\mathcal Av)=0.}
\]
In the continuation region \(v>g\), the last condition forces the PDE
equality. In the exercise region \(v=g\), the PDE residual need not vanish.
The terminal and spatial conditions are
\[
v(T,s)=(K-s)^+,\qquad v(t,0)=K,\qquad
\lim_{s\to\infty}v(t,s)=0.
\]
At zero stock price, immediate exercise pays \(K\); postponing that payment
cannot improve its discounted value when \(r\geq0\). On a numerical domain
\([0,S_{\max}]\), with \(S_{\max}>K\), replace the condition at infinity
by \(v(t,S_{\max})\approx0\), and check the truncation error by increasing
\(S_{\max}\).

\begin{warningbox}{An exercise grid is an approximation to continuous exercise}
For exercise restricted to \(0,h,2h,\ldots,T\), the Bermudan value satisfies
\[
v_h(t,s)=\max\{g(s),\E^\Q_{t,s}[e^{-rh}v_h(t+h,S_{t+h})]\}.
\]
For a finite \(h\), this forbids exercise between consecutive grid dates.
The continuous-time American value is recovered as the mesh tends to zero;
the same finite-step identity is not exact for unrestricted exercise.
\end{warningbox}

\begin{studentexercise}{identify the region}
For an American put, a numerical scheme reports \(g(s)=18\) and a
continuation estimate \(C=18.7\). Which region contains the node and what
is the option value?
\end{studentexercise}
\begin{shortsolution}
It is in the continuation region and \(v=18.7\). Intrinsic value is a lower
obstacle, not necessarily the price.
\end{shortsolution}

% ============================================================
\section{Computing American prices}
% ============================================================

\subsection{Finite differences: grid and explicit update}

Keep the no-dividend American put assumptions above. Choose
\[
s_i=i\Delta s,\quad \Delta s=S_{\max}/M,\qquad
t_n=n\Delta t,\quad \Delta t=T/N,
\]
and write \(V_i^n\approx v(t_n,s_i)\), \(g_i=(K-s_i)^+\).
Initialize \(V_i^N=g_i\), fix \(V_0^n=K\), \(V_M^n=0\), and march
backward from \(n=N-1\) to \(0\). At an interior stock node, use
\[
v_s(t_n,s_i)\approx\frac{V_{i+1}^n-V_{i-1}^n}{2\Delta s},\qquad
v_{ss}(t_n,s_i)\approx
\frac{V_{i+1}^n-2V_i^n+V_{i-1}^n}{(\Delta s)^2}.
\]
Let \(\mathcal A_h\) be the resulting discrete pricing operator.
An explicit continuation step evaluates it at the already known level:
\[
\frac{V_i^{n+1}-C_i^n}{\Delta t}
+(\mathcal A_hV^{n+1})_i=0.
\]
Since \(s_i=i\Delta s\), this gives the implementable update
\[
\boxed{
C_i^n=a_iV_{i-1}^{n+1}+b_iV_i^{n+1}+c_iV_{i+1}^{n+1},\qquad
V_i^n=\max\{g_i,C_i^n\},}
\]
where
\[
a_i=\frac{\Delta t}{2}(\sigma^2i^2-ri),\qquad
b_i=1-\Delta t(\sigma^2i^2+r),\qquad
c_i=\frac{\Delta t}{2}(\sigma^2i^2+ri).
\]
Thus each new value uses three neighbouring values at the next time level,
followed by the exercise test. Compute all interior nodes from the old
array before replacing it. The weights sum to \(1-r\Delta t\), the
first-order approximation to the one-step discount factor. At the end,
interpolate \(V^0\) to \(S_0\) if necessary. Nodes with
\(V_i^n=g_i>0\) identify the put exercise region; its largest stock node
approximates the free boundary.

\begin{warningbox}{Positivity and the explicit time-step restriction}
Require \(a_i,b_i,c_i\geq0\) for all \(1\leq i\leq M-1\). With
\(r\geq0\), sufficient conditions for the centered stencil are
\[
\sigma^2i^2\geq ri\quad\text{at every interior node},\qquad
\Delta t\leq\frac{1}{\sigma^2(M-1)^2+r}.
\]
Then the continuation update is monotone and its nonnegative weights sum
to at most one; taking the maximum with the payoff preserves monotonicity.
Reducing \(\Delta t\) fixes a negative \(b_i\), but cannot fix a negative
\(a_i\). If the drift dominates diffusion near zero, replace the centered
drift there by a monotone one-sided approximation and recompute the weights.
For \(r\geq0\), the forward difference
\(rs_i(V_{i+1}-V_i)/\Delta s\) gives nonnegative off-diagonal weights.
\end{warningbox}

\begin{boardexample}{one step before maturity}
Take \(K=100\), \(r=5\%\), \(\sigma=30\%\), \(\Delta s=5\),
\(\Delta t=1/252\), and \(S_{\max}=200\), so \(M=40\).
These parameters satisfy the positivity and time-step conditions.
Starting from the terminal payoff, two nodes give
\begin{center}
\begin{tabular}{c|c|c|c|l}
\toprule
\(s_i\) & \((V_{i-1}^N,V_i^N,V_{i+1}^N)\) & \(C_i^{N-1}\) & \(g_i\) & Action\\
\midrule
\(80\) & \((25,20,15)\) & \(19.9802\) & \(20\) & Exercise\\
\(100\) & \((5,0,0)\) & \(0.3472\) & \(0\) & Continue\\
\bottomrule
\end{tabular}
\end{center}
At \(s=80\), the payoff is locally linear and
\(C=20-rK\Delta t=19.9801587\). At \(s=100\),
\(a_{20}=\frac{1}{2\cdot252}(36-1)=0.0694444\), giving
\(C=5a_{20}=0.3472222\). These are numerical grid values, not exact
continuous-time prices.
\end{boardexample}

\subsection{Implicit Euler and discrete complementarity}

To avoid the explicit time-step restriction, evaluate spatial terms at the
unknown level \(n\). With the same \(a_i,c_i\), set
\[
d_i=1+\Delta t(\sigma^2i^2+r),\qquad
R_i^n=d_iV_i^n-a_iV_{i-1}^n-c_iV_{i+1}^n-V_i^{n+1}.
\]
The fully implicit obstacle problem is
\[
\boxed{V_i^n-g_i\geq0,\qquad R_i^n\geq0,\qquad
(V_i^n-g_i)R_i^n=0,\quad i=1,\ldots,M-1.}
\]
The endpoint values are fixed at \(K\) and \(0\). The residual is
\(-\Delta t\) times the discrete PDE residual, explaining its sign.
Provided \(a_i,c_i\geq0\), the tridiagonal system has positive diagonal,
nonpositive off-diagonals and strict diagonal dominance:
\(d_i-a_i-c_i=1+r\Delta t>0\). The explicit restriction on \(\Delta t\)
is no longer needed for monotonicity, although accuracy still requires
time and space refinement.

\begin{keyidea}{An implementable projected Gauss--Seidel solve}
At each time level, initialize \(U_i^{(0)}=V_i^{n+1}\), keep
\(U_0^{(k)}=K\), \(U_M^{(k)}=0\), and sweep \(i=1,\ldots,M-1\):
\[
U_i^{(k+1)}=\max\left\{g_i,
\frac{V_i^{n+1}+a_iU_{i-1}^{(k+1)}+c_iU_{i+1}^{(k)}}{d_i}
\right\}.
\]
Repeat the sweeps until
\[
\max_{1\leq i\leq M-1}
\left|\min\{U_i^{(k)}-g_i,\ R_i(U^{(k)})\}\right|\leq\varepsilon,
\]
then accept \(V^n=U^{(k)}\) and move to the preceding time level.
Here \(R_i(U)=d_iU_i-a_iU_{i-1}-c_iU_{i+1}-V_i^{n+1}\).
\end{keyidea}

One unconstrained implicit linear solve followed by a single pointwise
maximum does not generally solve this coupled complementarity problem:
changing one node changes its neighbours' PDE residuals. That procedure is
a splitting approximation; projected iteration enforces both conditions
together. Policy iteration and penalty methods provide other solvers.

{\footnotesize Further reading:
P.~Forsyth, \href{https://student.cs.uwaterloo.ca/~cs774/as2_09.pdf}{\emph{Numerical PDE Option Pricing}, CS~774, Assignment~2 (2009)},
for the complementarity formulation and implicit treatment of the obstacle.\par}

\subsection{Trees and least-squares Monte Carlo}

A recombining tree applies the Bellman maximum at each node and makes the
exercise region visible. Finite-difference schemes solve a discrete obstacle
problem: each grid value must satisfy both a PDE relation and \(v\geq g\).
Projection, penalty, and policy-iteration methods enforce this constraint.

For many risk factors, Longstaff--Schwartz approximates continuation values
from simulated risk-neutral paths:
\begin{enumerate}
  \item simulate paths at a set of exercise dates;
  \item set terminal cash flows to \(g(S_T)\);
  \item work backward and regress discounted future cash flows on basis
        functions of the current state, using in-the-money paths;
  \item exercise when \(g(S_{t_n})\) exceeds the fitted continuation value;
  \item average the discounted cash flows induced by the estimated policy.
\end{enumerate}

\begin{keyidea}{What LSM estimates}
Regression estimates the conditional continuation value
\[
C_n(s)\approx\E^\Q[e^{-r(t_{n+1}-t_n)}V_{n+1}\mid S_{t_n}=s],
\]
not the price directly. Poor basis functions, sparse exercise dates, or too
few paths can produce a poor stopping policy.
\end{keyidea}

\begin{warningbox}{Mandatory numerical checks}
Refine the time grid and verify
\[
V^{\rm Am}\geq V^{\rm Eur}\geq0,\qquad
V^{\rm Am,put}\geq(K-S_0)^+.
\]
For a non-dividend-paying call, the American method should converge to the
European benchmark. A material gap is a diagnostic warning.
\end{warningbox}

\begin{studentexercise}{diagnose two numerical errors}
You price an up-and-in American put by Monte Carlo with daily monitoring.
Name one issue caused by the barrier and one caused by the American feature.
\end{studentexercise}
\begin{shortsolution}
A path may cross the barrier between monitoring dates, creating monitoring
bias; refine the grid or use a bridge correction. The American feature
requires a stopping policy, so test exercise dates, regression bases, and
out-of-sample policy performance.
\end{shortsolution}

% ============================================================
\section{Synthesis and guided practice}
% ============================================================

\begin{keyidea}{A practical workflow}
\begin{enumerate}
  \item Write the payoff and permitted exercise dates.
  \item Compare exercise and risk-neutral continuation values.
  \item In a tree, apply the Bellman maximum from maturity backward.
  \item In continuous time, locate stopping and continuation regions.
  \item Check numerical output against bounds and refinements.
\end{enumerate}
\end{keyidea}

\Needspace{11\baselineskip}
\begin{studentexercise}{final stopping-policy check}
An American put has immediate exercise value \(12\) and continuation value
\(11.5\). A colleague adds them and reports an option value of \(23.5\).
Explain the error, state the correct value, and give the action.
\end{studentexercise}
\begin{shortsolution}
Exercise and continuation are mutually exclusive policies, so the recursion
takes their maximum:
\[
V=\max\{12,11.5\}=12.
\]
The holder exercises. The continuation value belongs to a different policy.
\end{shortsolution}

\begin{keyidea}{Final synthesis}
\[
\boxed{\text{American pricing}
=\text{no-arbitrage valuation}+\text{an optimal stopping decision}.}
\]
Dynamic programming, the Snell envelope, and the free-boundary PDE are
three compatible ways to represent the same exercise right.
\end{keyidea}
\end{document}
